% !TEX program = xelatex

\documentclass[14pt,a4paper]{extreport}
\usepackage{timcw}
\usepackage{fontspec}
\usepackage{polyglossia}
\setmainfont{Times New Roman}[Ligatures=TeX]
\setsansfont{Arial}[Ligatures=TeX]
\setmonofont{Courier New}[Ligatures=TeX]
\setmainlanguage{russian}
\setotherlanguage{english}
\begin{document}
\begin{singlespacing}
% Титульный лист:
% #1 – тема, #2 – группа, #3 – ФИО, #4 – город и год
\cwtitle{Разработка веб-приложения для проведения викторин с геймификацией}{Д-Э-311}{Конюхов М.В.}{Москва 2026}{09.03.02 Информационные системы и технологии}{Компьютерные науки и интеллектуальный анализ данных}
\newpage

\renewcommand{\contentsname}{\hfill Содержание\hfill}

\tableofcontents
\end{singlespacing}
\newpage

\intro

В связи с быстрым развитием информационных технологий и широким использованием интернета веб-приложения стали главным средством автоматизации учебных и рабочих процессов. Одним из многообещающих направлений является создание систем проверки знаний, в которых используются элементы геймификации. Такой подход повышает заинтересованность пользователей благодаря игровым элементам: начислению баллов, получению уровней, выдаче наград и ведению списков лучших участников.

Актуальность работы вызвана растущей потребностью в интерактивных платформах для тестирования знаний в образовании и корпоративном обучении. Обычные способы проверки знаний часто не дают достаточной мотивации участникам, тогда как игровые системы показывают значительно лучшие результаты по удержанию пользователей и качеству усвоения учебного материала.

Цель курсового проекта заключается в создании веб-приложения для проведения викторин с игровыми элементами на основе современного набора технологий.

Для выполнения поставленной цели требуется решить следующие задачи:

\begin{itemize}
    \item Изучить правила построения информационных систем и способы создания веб-приложений;
    \item Проанализировать существующие решения в области платформ для тестирования и геймификации;
    \item Обосновать выбор инструментов и технологий, используемых при разработке;
    \item Разработать архитектуру системы, схему базы данных и пользовательский интерфейс;
    \item Создать веб-приложение с системой входа пользователей, модулем викторин и игровыми механиками;
    \item Проверить работоспособность разработанной системы и проанализировать полученные результаты.
\end{itemize}

Объектом исследования является процесс проектирования и создания информационных систем на основе веб-технологий.

Предметом исследования выступают методы и средства построения веб-приложения для викторин с элементами геймификации.
\newpage
\chapter[Глава \thechapter. Теоретическая часть]{Теоретическая часть}

\section{Обзор предметной области}
Веб-викторины представляют собой интерактивные информационные системы, предназначенные для проверки знаний, формирования навыков и повышения вовлечённости пользователей. Такие приложения могут функционировать в двух режимах:
\begin{itemize}
    \item синхронном, когда несколько участников одновременно выполняют задания под управлением ведущего;
    \item асинхронном, где пользователь проходит задания в удобное для себя время.
\end{itemize}

Формат вопросов варьируется от простых вариантов множественного выбора до открытых вопросов и заданий на заполнение пропусков. Выбор формата влияет на способы оценивания и анализа полученных результатов.

В научной литературе геймификация рассматривается как инструмент, который при правильном использовании способен повышать внутреннюю мотивацию пользователей, обеспечивать быструю обратную связь и поддерживать эффект «потока» -- то есть оптимальное соотношение между сложностью задания и уровнем навыков участника.

Области применения веб-викторин разнообразны -- от образовательных и корпоративных учреждений, до публичных мероприятий и маркетинговых кампаний.

В каждом из этих контекстов требования к функциям системы и показателям эффективности различаются. Для образовательных задач важна точность оценивания и анализ учебных достижений, тогда как для развлекательных мероприятий ключевыми являются вовлечённость аудитории и устойчивый интерес к платформе.

Основные требования предметной области включают:
\begin{itemize}
    \item поддержку различных типов вопросов;
    \item надёжный учёт результатов;
    \item механизмы формирования рейтингов и отчётов;
    \item возможность масштабируемого проведения сессий в реальном времени.
\end{itemize}

Также важное значение имеют вопросы безопасности и конфиденциальности данных пользователей, особенно в образовательной и корпоративной среде. Для оценки эффективности систем применяются такие метрики, как вовлечённость, частота повторного использования, доля успешно пройденных заданий и качество обучения. Это позволяет соотнести выбранные игровые механики с реальными учебными или бизнес-целями.

\section{Теоретические основы геймификации}
Геймификация определяется как применение игровых элементов и приёмов в неигровых ситуациях с целью повышения мотивации и вовлечённости пользователей. В основе её эффективности лежат психологические и педагогические теории мотивации. Теория самоопределения выделяет три главные потребности человека:
\begin{itemize}
    \item автономия -- возможность самостоятельно выбирать действия;
    \item компетентность -- ощущение успешности и умения справляться с задачами;
    \item связанность -- чувство принадлежности к группе и взаимодействие с другими.
\end{itemize}

Удовлетворение этих потребностей способствует росту внутренней мотивации (Ryan \& Deci, 2000). Концепция «потока» показывает, что наилучший опыт достигается при равновесии между сложностью задания и уровнем навыков пользователя (Csikszentmihalyi, 1990). Бихевиористские модели делают упор на роль подкрепления: поощрения усиливают желаемое поведение и делают его более частым.

Игровые механики выполняют три взаимосвязанные функции: дают обратную связь, выстраивают структуру прогресса, создают социальное взаимодействие.

В контексте викторин можно воспользоваться следующими элементами:
\begin{itemize}
    \item очки за правильные ответы;
    \item бонусы за быстроту реакции;
    \item бонусы за серии верных ответов;
    \item уровни пользователей;
    \item достижения (награды).
\end{itemize}

Такие механики помогут пользователю ощущать прогресс и рост своей компетентности, побудят возвращаться к заданиям и повысят активность во время прохождения викторины. Например, система очков и бонусов за быстрый ответ закрепляет оперативность и быстроту мышления, тогда как уровни и достижения поддерживают долгосрочную мотивацию, показывая шаги учебного или профессионального развития.

При разработке игровых механик важно учитывать их влияние на обучение. Механики должны усиливать учебную цель, а не заменять её. Очки и таблицы лидеров могут повышать вовлечённость, но если делать слишком большой упор на соревновательность, есть риск, что пользователи будут хуже усваивать материал. Серии правильных ответов и ежедневные задания помогают выработать привычку и регулярно заниматься, что положительно сказывается на запоминании информации. Одновременно необходимо давать обратную связь об ошибках: краткое пояснение после ответа позволяет превратить поведенческую реакцию в осознанное знание.

Правила геймификации должны быть чёткими и понятными. Пользователю необходимо ясно понимать:
\begin{itemize}
    \item за что начисляются очки;
    \item как рассчитываются бонусы;
    \item какие условия нужно выполнить для получения достижения или повышения уровня.
\end{itemize}

Баланс наград является очень важным условием. Если давать слишком много очков за простые действия, ценность достижений снижается и мотивация падает. Если награды слишком редкие или непонятные, они не будут побуждать к действию.

Социальные аспекты геймификации сильно влияют на вовлечённость. Таблицы лидеров и групповая динамика усиливают соревновательные и кооперативные мотивы. Однако нужно учитывать, что разные пользователи по-разному реагируют на конкуренцию. Для одних людей открытое соревнование служит мотиватором. Для других оно, наоборот, снижает желание участвовать.

Поэтому полезно будет предусмотреть настройки приватности, а также механики, ориентированные на личный прогресс, а не только на сравнение с другими участниками.

Оценка эффективности геймификации должна опираться на набор показателей, связанных как с вовлечённостью, так и с учебными результатами. Ключевыми показателями являются: частота и длительность сессий, количество повторных обращений к платформе, доля завершённых викторин, динамика среднего балла, сохранение знаний с течением времени.

Анализ этих данных позволяет проверить, какие игровые механики дают наибольший эффект, и при необходимости скорректировать систему наград с учётом выявленных моделей поведения.

\section{Анализ существующих решений}

Анализ существующих платформ для проведения викторин позволяет выявить типичные подходы к устройству интерфейса, игровым механикам и сбору статистики, а также определить часто встречающиеся ограничения. Это обосновывает необходимость создания нового решения. Ниже рассмотрены наиболее известные сервисы с точки зрения их возможностей, технических особенностей, достоинств и недостатков.

\paragraph{Kahoot!}
Данный сервис широко известен как инструмент для синхронных викторин, проводимых в реальном времени. Интерфейс ориентирован на простоту и вовлечение участников: ведущий запускает сессию, игроки подключаются по коду и отвечают на вопросы одновременно. Геймификация реализована через:
\begin{itemize}
    \item начисление очков за правильные ответы и быстрое реагирование;
    \item таблицы лидеров в рамках конкретной сессии.
\end{itemize}
\textit{Сильные стороны Kahoot!:}
\begin{itemize}
    \item низкий порог входа для пользователей;
    \item яркий визуальный дизайн;
    \item высокая вовлечённость на очных или дистанционных занятиях.
\end{itemize}
\textit{Недостатки:}
\begin{itemize}
    \item ограниченная аналитика по долгосрочным результатам;
    \item слабая поддержка сложных типов вопросов;
    \item зависимость игровой динамики от синхронного участия всех игроков.
\end{itemize}

\paragraph{Quizizz}
Эта платформа сочетает в себе возможности синхронных и асинхронных викторин. Quizizz поддерживает:
\begin{itemize}
    \item домашние задания;
    \item автоматическое оценивание;
    \item более развитые отчёты об успеваемости по сравнению с сервисами, ориентированными только на живые игры.
\end{itemize}
Встроенные игровые элементы включают очки, мемы, аватары и таблицы лидеров. Quizizz даёт больше гибкости в настройке прохождения и распределении заданий. Однако настройка логики начисления наград и сложных достижений остаётся ограниченной. Для задач строгого оценивания и детальной аналитики это может быть недостатком.

\paragraph{Moodle (модуль Quiz)}
Moodle представляет собой более профессиональное решение для учебных заведений. Ключевые преимущества:
\begin{itemize}
    \item тесная интеграция с системой управления обучением (LMS);
    \item гибкие типы вопросов;
    \item масштабируемые механизмы оценивания;
    \item возможность расширения через плагины.
\end{itemize}
Moodle обеспечивает детальную статистику по вопросам и пользователям, а также удобные инструменты для модерации. Вместе с тем:
\begin{itemize}
    \item пользовательский интерфейс часто воспринимается как громоздкий;
    \item настройка и администрирование требуют значительных усилий;
    \item возможности геймификации в базовой версии ограничены и обычно добавляются через сторонние расширения.
\end{itemize}

\paragraph{Socrative и аналогичные системы}
Socrative и другие системы оперативного сбора ответов в классе ориентированы на:
\begin{itemize}
    \item быстрый сбор ответов во время занятий;
    \item простоту использования для преподавателя.
\end{itemize}
Они удобны для быстрой проверки понимания материала, но не рассчитаны на долгосрочное накопление достижений или сложную систему уровней.

\paragraph{Mentimeter}
Данный сервис больше нацелен на интерактивные презентации и опросы. Функционал для оценки знаний и развитых игровых механик в нём ограничен.

\paragraph{Google Forms и универсальные формы}
Эти инструменты служат простым способом создания тестов и опросов. Их преимущества:
\begin{itemize}
    \item доступность;
    \item простота использования;
    \item интеграция с экосистемой Google.
\end{itemize}
Однако они практически не приспособлены для синхронных игр, не поддерживают таблицы лидеров в реальном времени и не имеют встроенной геймификации. Аналитика ограничена базовыми отчётами и требует дополнительной обработки данных.

\paragraph{Технические особенности коммерческих решений}
Многие коммерческие решения используют:
\begin{itemize}
    \item web-сокеты или собственные сервисы реального времени для синхронных сессий;
    \item облачные хранилища для медиафайлов;
    \item внешние сервисы аналитики.
\end{itemize}
Платная модель и ограниченная открытость API в некоторых продуктах создают препятствия для интеграции с системами обучения и корпоративными системами. Вопросы приватности и хранения персональных данных также различаются:
\begin{itemize}
    \item у образовательных платформ часто есть готовые механизмы для соблюдения местных законов;
    \item у развлекательных сервисов политика в отношении данных может быть менее строгой.
\end{itemize}

\paragraph{Выявленные пробелы}
Проведённый обзор позволяет выделить несколько направлений для создания нового приложения:
\begin{enumerate}
    \item Не хватает гибкой и понятной системы геймификации, которая объединяла бы как сессионные механики (очки, бонусы за скорость, таблицы лидеров), так и долгосрочные стимулы (уровни, достижения, ежедневные серии).
    \item Многие решения предлагают либо мощную аналитику без удобной геймификации, либо, наоборот, сильную вовлечённость без глубоких инструментов оценки и отслеживания прогресса.
    \item Ограниченный набор типов вопросов, слабая настройка сложности и недостаток адаптивных сценариев затрудняют использование платформ там, где важна индивидуальная настройка обучения.
    \item Отсутствие открытых API и сложность подключения к системам обучения или единому входу (SSO) мешают внедрению решений в существующую образовательную инфраструктуру.
\end{enumerate}

\paragraph{Место разрабатываемого решения}
В рамках данного проекта реализованы следующие механики:
\begin{itemize}
    \item очки за правильные ответы;
    \item бонусы за быстрое реагирование;
    \item бонусы за серии правильных ответов;
    \item уровни пользователей;
    \item достижения;
    \item ежедневные серии.
\end{itemize}
Эти решения направлены именно на заполнение указанных пробелов. Сочетание синхронного и асинхронного режимов, детализированной аналитики попыток и гибкой настройки правил начисления наград создаёт ценность, которая отличает данную разработку от популярных аналогов. Кроме того, выбранные архитектурные решения обеспечивают прозрачность правил, возможность экспорта данных и интеграцию с существующими образовательными системами, что повышает практическую пользу от разработки.


\newpage
\chapter[Глава \thechapter. Практическая часть]{Практическая часть}
text (\ref{summa})

\section{Пункт}
text (\ref{alphabeta})

\section{Пункт}

\newpage
\conclusion

Текст заключения.

\newpage
\centeredrefname
\clearpage
\addcontentsline{toc}{chapter}{Список использованных источников}
\begin{thebibliography}{99}
\bibitem{book3} Author Book3
\bibitem{book1} Author Book1
\bibitem{book2} Author Book2
\end{thebibliography}

\newpage
\applications

Текст приложений.

\end{document}